%============================================================
% DANH MỤC THUẬT NGỮ VÀ TỪ VIẾT TẮT
%============================================================
\chapter*{DANH MỤC THUẬT NGỮ VÀ TỪ VIẾT TẮT}
\addcontentsline{toc}{chapter}{Danh mục thuật ngữ và từ viết tắt}

\begin{longtable}{p{3cm} p{11cm}}
\toprule
\textbf{Thuật ngữ} & \textbf{Ý nghĩa} \\
\midrule
\endfirsthead
\toprule
\textbf{Thuật ngữ} & \textbf{Ý nghĩa} \\
\midrule
\endhead
\bottomrule
\endfoot

AI      & Trí tuệ nhân tạo (Artificial Intelligence), mô phỏng các quá
          trình suy nghĩ và học tập của con người bằng máy móc \\[4pt]

AIoT    & Trí tuệ nhân tạo kết hợp Internet vạn vật (Artificial Intelligence
          of Things), tích hợp AI vào các thiết bị IoT để xử lý và phân
          tích dữ liệu tại chỗ \\[4pt]

AQI     & Chỉ số chất lượng không khí (Air Quality Index), thang đo mức
          độ ô nhiễm không khí ảnh hưởng đến sức khỏe con người \\[4pt]

API     & Giao diện lập trình ứng dụng (Application Programming Interface),
          tập hợp các giao thức cho phép các phần mềm giao tiếp với nhau \\[4pt]

BLE     & Công nghệ truyền thông không dây tầm ngắn tiêu thụ năng lượng
          thấp (Bluetooth Low Energy) \\[4pt]

CNN     & Mạng nơ-ron tích chập (Convolutional Neural Network), loại mạng
          nơ-ron thường dùng cho xử lý ảnh và tín hiệu tuần tự \\[4pt]

CO\textsubscript{2}
        & Khí carbon dioxide, một trong các chất ô nhiễm không khí phổ biến \\[4pt]

DHT22   & Cảm biến nhiệt độ và độ ẩm kỹ thuật số độ chính xác cao \\[4pt]

Edge Computing
        & Mô hình tính toán phân tán, xử lý dữ liệu gần nguồn phát sinh
          thay vì tập trung tại trung tâm dữ liệu \\[4pt]

LSTM    & Bộ nhớ ngắn hạn dài (Long Short-Term Memory), loại mạng nơ-ron
          hồi tiếp chuyên xử lý dữ liệu chuỗi thời gian \\[4pt]

IoT     & Internet vạn vật (Internet of Things), mạng lưới các thiết bị
          vật lý được kết nối internet để thu thập và chia sẻ dữ liệu \\[4pt]

MAE     & Sai số tuyệt đối trung bình (Mean Absolute Error), độ đo đánh
          giá hiệu năng mô hình dự báo \\[4pt]

ML      & Học máy (Machine Learning), lĩnh vực AI cho phép máy tính học
          từ dữ liệu mà không cần lập trình cụ thể \\[4pt]

MQTT    & Giao thức truyền tin nhẹ (Message Queuing Telemetry Transport),
          thường dùng trong các hệ thống IoT \\[4pt]

MQ-135  & Cảm biến khí đa năng, phát hiện NH\textsubscript{3}, NO\textsubscript{x},
          cồn, benzen, khói và CO\textsubscript{2} \\[4pt]

PM2.5   & Hạt bụi mịn có đường kính nhỏ hơn hoặc bằng 2{,}5 micromet,
          gây hại nghiêm trọng cho đường hô hấp \\[4pt]

PM10    & Hạt bụi có đường kính nhỏ hơn hoặc bằng 10 micromet \\[4pt]

PMS5003 & Cảm biến bụi mịn laser đo PM1.0, PM2.5 và PM10 \\[4pt]

RMSE    & Căn bậc hai của sai số bình phương trung bình (Root Mean Squared
          Error), độ đo đánh giá hiệu năng mô hình dự báo \\[4pt]

TFLite  & TensorFlow Lite, framework chạy mô hình ML trên thiết bị nhúng
          và di động với tài nguyên hạn chế \\[4pt]

TinyML  & Triển khai các mô hình học máy nhỏ gọn trực tiếp trên vi điều
          khiển có tài nguyên hạn chế \\[4pt]

\end{longtable}
